\chapter{Fracciones: comparar y sumar}\label{ch:g7:fractions}

Grado 6 introducido \omterm{def:g6:fractions:def}{fracciones} como acciones y como exacto \omterm{def:g3:division:remainder}{cocientes}
(\cref{ch:g6:fractions}). Este capítulo aprende a \emph{calcular} con
ellos: reconocer la igualdad \omterm{def:g6:fractions:def}{fracciones}, comparando, sumando y restando
--- por ahora con amigable \omterm{def:g4:fractions:def}{denominadores}; el caso general entra
\cref{ch:g9:fractions}.

\section{fracciones iguales}

\begin{theorem}[fracciones iguales]\label{thm:g7:fractions:equal}
A \omterm{def:g6:fractions:def}{fracción} no cambia cuando su \omterm{def:g4:fractions:def}{numerador} y \omterm{def:g4:fractions:def}{denominador} son
multiplicado o dividido por el mismo número distinto de cero:
\[
\frac{a}{b} = \frac{a \times k}{b \times k} .
\]
\end{theorem}
\admitted

\begin{example}\label{ex:g7:fractions:equal}
$\dfrac{3}{4} = \dfrac{3 \times 5}{4 \times 5} = \dfrac{15}{20}$, y
en la otra dirección $\dfrac{42}{30} = \dfrac{42 \div 6}{30 \div 6}
= \dfrac{7}{5}$ (simplificado por $6$).

Para probar si $\dfrac{8}{12}$ y $\dfrac{10}{15}$ son iguales,
simplificar ambos: $\dfrac{8}{12} = \dfrac{2}{3}$ y
$\dfrac{10}{15} = \dfrac{2}{3}$ --- si, igual.
\end{example}

\section{Comparar fracciones}

\begin{method}[Comparando dos fracciones]\label{met:g7:fractions:compare}
\begin{enumerate}
  \item \emph{Mismo \omterm{def:g4:fractions:def}{denominador}}: comparar el \omterm{def:g4:fractions:def}{numeradores}:
        $\frac{5}{9} < \frac{7}{9}$.
  \item \emph{Uno \omterm{def:g4:fractions:def}{denominador} es un \omterm{def:g5:division:multiple}{múltiple} del otro}: reescribir el
        más grueso \omterm{def:g6:fractions:def}{fracción}, luego compara:
        $\frac{5}{6}$ vs $\frac{7}{12}$: $\frac56 = \frac{10}{12} >
        \frac{7}{12}$.
  \item \emph{Comparar con $1$ o $\frac12$} cuando sea posible:
        $\frac{9}{8} > 1 > \frac{7}{9}$ se asienta
        $\frac98$ vs $\frac79$ instantáneamente.
\end{enumerate}
\end{method}

\begin{omfigure}
\begin{tikzpicture}[scale=0.62]
  \foreach \i in {0,1,2,3,4,5} \draw[omExo] (\i*2,0) rectangle (\i*2+2,1.1);
  \foreach \i in {0,1,2,3,4} \fill[omDef!35] (\i*2+0.05,0.05)
    rectangle (\i*2+1.95,1.05);
  \draw[thick] (0,0) rectangle (12,1.1);
  \node[right, font=\small] at (12.3,0.55) {$\frac56$};
  \begin{scope}[yshift=-1.9cm]
  \foreach \i in {0,1,2,3,4,5,6,7,8,9,10,11} \draw[omExo] (\i,0) rectangle (\i+1,1.1);
  \foreach \i in {0,1,2,3,4,5,6} \fill[omThm!35] (\i+0.05,0.05)
    rectangle (\i+0.95,1.05);
  \draw[thick] (0,0) rectangle (12,1.1);
  \node[right, font=\small] at (12.3,0.55) {$\frac{7}{12}$};
  \end{scope}
\end{tikzpicture}

{\small Comparando $\frac56$ y $\frac{7}{12}$ en barras gemelas: corte
los sextos en \omterm{def:g3:division:half}{medio} muestra $\frac56 = \frac{10}{12}$, más que
$\frac{7}{12}$.}
\end{omfigure}

\section{Sumar y restar}

\begin{theorem}[Mismo denominador]\label{thm:g7:fractions:add}
\omterm{def:g6:fractions:def}{fracciones} con el mismo \omterm{def:g4:fractions:def}{denominador} se suman (o restan) por
sumando (o restando) el \omterm{def:g4:fractions:def}{numeradores}:
\[
\frac{a}{d} + \frac{b}{d} = \frac{a + b}{d},
\qquad
\frac{a}{d} - \frac{b}{d} = \frac{a - b}{d} .
\]
\end{theorem}

\begin{proof}[Por qué]
$ a $ partes de tamaño $\frac1d $ más $ b $ partes del mismo tamaño hacen
$ a + b $ partes de ese tamaño.
\end{proof}

\begin{method}[Diferentes denominadores]\label{met:g7:fractions:add}
cuando uno \omterm{def:g4:fractions:def}{denominador} es un \omterm{def:g5:division:multiple}{múltiple} del otro:
\begin{enumerate}
  \item reescribir el \omterm{def:g6:fractions:def}{fracción} con el mas pequeño \omterm{def:g4:fractions:def}{denominador} entonces ambos tienen
        el más grande (\cref{thm:g7:fractions:equal});
  \item sumar o restar el \omterm{def:g4:fractions:def}{numeradores};
  \item Simplifique el resultado si es posible.
\end{enumerate}
\end{method}

\begin{example}\label{ex:g7:fractions:add}
Calcular $\dfrac{3}{4} + \dfrac{5}{8}$, paso a paso:
\begin{enumerate}
  \item $8$ es un \omterm{def:g5:division:multiple}{múltiple} de $4$: reescribir
        $\dfrac34 = \dfrac{3 \times 2}{4 \times 2} = \dfrac{6}{8}$;
  \item agregar: $\dfrac{6}{8} + \dfrac{5}{8} = \dfrac{11}{8}$;
  \item nada que simplificar: el resultado es $\dfrac{11}{8}$, es decir.\
        $1 + \dfrac38$.
\end{enumerate}
Otro: $2 - \dfrac{4}{5} = \dfrac{10}{5} - \dfrac{4}{5} =
\dfrac{6}{5}$ (escribe el número entero sobre el \omterm{def:g4:fractions:def}{denominador} $5$
primero).
\end{example}

\begin{example}[Una trampa clásica]\label{ex:g7:fractions:trap}
$\dfrac12 + \dfrac13$ es \emph{no} $\dfrac{2}{5}$! Añadiendo \omterm{def:g4:fractions:def}{numeradores}
y \omterm{def:g4:fractions:def}{denominadores} por separado está mal --- la respuesta $\frac25$ sería
más pequeño que $\frac12$, absurdo cuando \emph{agregando} algo positivo
a $\frac12$. (El correcto \omterm{def:g1:addition:def}{suma}, $\frac56$, necesita un común \omterm{def:g4:fractions:def}{denominador}
$6$; el método general está en \cref{ch:g9:fractions}.)
\end{example}

\section{Multiplicar una fracción por un número}

\begin{proposition}[Fracción multiplicada por un número]\label{prop:g7:fractions:mult}
por un número $ k $:
\[
k \times \frac{a}{b} = \frac{k \times a}{b} .
\]
Tomando $\frac ab $ \emph{de} una cantidad significa multiplicar por
$\frac ab $.
\end{proposition}
\admitted

\begin{example}\label{ex:g7:fractions:mult}
$6 \times \dfrac{2}{3} = \dfrac{12}{3} = 4$: seis por dos tercios es
cuatro enteros. y ``$\frac34$ de $60$ euros'' es
$\frac34 \times 60 = \frac{180}{4} = 45$ euros --- la misma respuesta que
dividiendo por $4$ luego multiplicando por $3$
(\cref{met:g6:fractions:ofquantity}).
\end{example}

\section{Ejercicios}

\begin{exercise}[$\star $]\label{exo:g7:fractions:1}
Completo: $\dfrac{2}{5} = \dfrac{?}{20}$; \quad
$\dfrac{18}{24} = \dfrac{3}{?}$; \quad
$\dfrac{7}{3} = \dfrac{28}{?}$; \quad
$\dfrac{?}{9} = \dfrac{20}{36}$.
\end{exercise}

\begin{exercise}[$\star $]\label{exo:g7:fractions:2}
Simplifique tanto como sea posible:
$\dfrac{12}{16}$; \quad $\dfrac{30}{45}$; \quad $\dfrac{27}{36}$;
\quad $\dfrac{48}{12}$.
\end{exercise}

\begin{exercise}[$\star $]\label{exo:g7:fractions:3}
Son $\dfrac{15}{35}$ y $\dfrac{21}{49}$ ¿igual? Y $\dfrac{16}{28}$
y $\dfrac{20}{36}$? Justifica simplificando.
\end{exercise}

\begin{exercise}[$\star $]\label{exo:g7:fractions:4}
Compara (escribe el razonamiento, no solo la respuesta):
\[
\frac{7}{11} \text{ y } \frac{9}{11}; \qquad
\frac{5}{6} \text{ y } \frac{11}{18}; \qquad
\frac{13}{12} \text{ y } \frac{19}{20}.
\]
\end{exercise}

\begin{exercise}[$\star $]\label{exo:g7:fractions:5}
Calcula y simplifica si es posible:
\[
\frac{5}{9} + \frac{2}{9}, \qquad
\frac{11}{7} - \frac{4}{7}, \qquad
\frac{5}{8} + \frac{7}{8} .
\]
\end{exercise}

\begin{exercise}[$\star $]\label{exo:g7:fractions:6}
Calcular usando \cref{met:g7:fractions:add}:
\[
\frac{2}{3} + \frac{5}{6}, \qquad
\frac{7}{10} - \frac{2}{5}, \qquad
\frac{3}{4} + \frac{5}{12} .
\]
\end{exercise}

\begin{exercise}[$\star $]\label{exo:g7:fractions:7}
Calcular:
\[
3 - \frac{5}{4}, \qquad
1 + \frac{3}{8}, \qquad
2 - \frac{7}{6} .
\]
\end{exercise}

\begin{exercise}[$\star $]\label{exo:g7:fractions:8}
Calcular:
\[
5 \times \frac{3}{10}, \qquad
\frac{7}{8} \times 4, \qquad
\frac{2}{3} \text{ of } 45 .
\]
\end{exercise}

\begin{exercise}[$\star\star $]\label{exo:g7:fractions:9}
marc comió $\frac14$ de una pizza y julie $\frac38$ de la misma pizza.
Qué \omterm{def:g6:fractions:def}{fracción} ¿Qué cantidad de pizza comieron juntos? Qué \omterm{def:g6:fractions:def}{fracción} es
¿Izquierda?
\end{exercise}

\begin{exercise}[$\star\star $]\label{exo:g7:fractions:10}
una botella contiene $\frac34$ L de jugo. Lea se derrama $\frac16$ l
dos veces. ¿Cuánto jugo queda? (Común \omterm{def:g4:fractions:def}{denominador}: $12$.)
\end{exercise}

\begin{exercise}[$\star\star $]\label{exo:g7:fractions:11}
Explique, con el argumento de \cref{ex:g7:fractions:trap}, por qué
$\frac35 + \frac12$ no puede ser $\frac{4}{7}$, sin calcular el
correcto \omterm{def:g1:addition:def}{suma}.
\end{exercise}

\begin{exercise}[$\star\star\star $]\label{exo:g7:fractions:12}
Calcular el \omterm{def:g1:addition:def}{suma}
\[
\frac12 + \frac14 + \frac18 + \frac{1}{16}
\]
paso a paso (de izquierda a derecha). Observa el patrón del parcial. \omterm{def:g1:addition:def}{sumas}:
¿A qué número se acercan cuando hay más y más términos?
$\frac{1}{32}, \frac{1}{64}, \dots $ se agregan?
\end{exercise}

\section{Problema: Las fracciones escondidas en un compás de música.}

\begin{problem}[{Problema de fin de semana --- los valores de las notas son fracciones que deben sumar exactamente, el conteo de ritmos oculta una secuencia famosa y la paridad gobierna el ritmo de los Balcanes}]
\label{pb:g7:fractions:1}
La música está escrita en \omterm{def:g6:fractions:def}{fracciones}. A \emph{nota completa} dura $1$
medida (en el tiempo común "cuatro-cuatro"); a \emph{\omterm{def:g3:division:half}{medio} nota}
dura $\frac12$; entonces viene el \emph{\omterm{def:g3:division:half}{cuarto} nota} $\frac14$, el
\emph{octava nota} $\frac18$ y el \emph{semicorchea}
$\frac{1}{16}$. Una ley inviolable: \emph{las duraciones dentro de un
La medida debe sumar exactamente el total de la medida.} --- $1$ en
cuatro por cuatro veces. Cada pregunta a continuación es la aritmética de este capítulo.
(\cref{thm:g7:fractions:add}, \cref{met:g7:fractions:add}) establecido en
música; en el camino te encontrarás con una secuencia de conteo descubierta por
Indian scholars a thousand years ago.

\medskip
\noindent\textbf{Parte I --- Notes that must add up.}
\begin{enumerate}
  \item Comprueba que cada una de estas medidas es legal en cuatro-cuatro
        time: (a) \omterm{def:g3:division:half}{medio} $+$ \omterm{def:g3:division:half}{cuarto} $+$ \omterm{def:g3:division:half}{cuarto}; (b) \omterm{def:g3:division:half}{cuarto} $+$
        octavo $+$ octavo $+$ \omterm{def:g3:division:half}{medio}.
  \item Un copista escribió: \omterm{def:g3:division:half}{medio} $+$ \omterm{def:g3:division:half}{cuarto} $+$ octavo. Qué
        \omterm{def:g6:fractions:def}{fracción} del compás está escrito, y cuál es el único
        note completes it?
  \item ¿Cuántas semicorcheas ocupan un compás completo? como
        ¿Cuántas semicorcheas forman una corchea?
  \item A \emph{punto} después agrega una nota \omterm{def:g3:division:half}{medio} the note's value: a
        punteado \omterm{def:g3:division:half}{medio} es $\frac12 + \frac14$, un punteado \omterm{def:g3:division:half}{cuarto} es
        $\frac14 + \frac18$. Calcule ambos valores.
  \item Un vals se escribe en compás de tres por cuatro: cada compás debe
        total $\frac34$. Compruebe que un solo punto \omterm{def:g3:division:half}{medio} llena un
        medida de vals y anota otras dos medidas de vals legales.
        medidas usando solo \omterm{def:g3:division:half}{cuarteles} y octavos.
\end{enumerate}

\medskip
\noindent\textbf{Parte II --- The copyist's workshop.}
\begin{enumerate}[resume]
  \item Un compás cuatro por cuatro contiene: punteado \omterm{def:g3:division:half}{cuarto} $+$ octavo
        $+$ \omterm{def:g3:division:half}{cuarto}. ¿Qué está faltando?
  \item ¿Cuál dura más, una octava punteada o una \omterm{def:g3:division:half}{cuarto} ¿nota?
        Comparar mediante dieciseisavos.
  \item A \emph{atar} pega duraciones juntas en una larga
        sonido: un \omterm{def:g3:division:half}{medio} nota atada a una octava ¿cuanto dura? como
        ¿Cuántos dieciseisavos es eso?
  \item Las notas de un compás suman $\frac{11}{16}$; el
        \omterm{def:g3:division:remainder}{resto} es el silencio (un \emph{descansar}). ¿Cuánto dura el
        descansar?
  \item Un baterista llena uno \omterm{def:g3:division:half}{cuarto} nota ($\frac{4}{16}$ de la
        compás) usando sólo octavos ($\frac{2}{16}$) y
        dieciseisavos ($\frac{1}{16}$), orden importante: para
        instancia octava--decimosexta--decimosexta, o
        decimosexto – octavo – decimosexto. Compruebe que ambos ejemplos
        son legales, luego enumere \emph{todo} Los ritmos posibles.
        ¿Cuantos hay?
\end{enumerate}

\medskip
\noindent\textbf{Parte III --- A thousand-year-old sequence, and
an \omterm{def:g3:numbers:evenodd}{extraño} measure.}
\begin{enumerate}[resume]
  \item Contemos los ritmos del baterista para notas más largas.
        Cada ritmo termina con una semicorchea o con una
        octavo; explica por qué esto da la regla: (ritmos
        relleno $ n $ dieciseisavos) $=$ (ritmos llenando $ n - 1$)
        $+$ (ritmos llenando $ n - 2$). A partir de $1$ ritmo
        por un decimosexto y $2$ para dos, construye la mesa hasta
        $ n = 8$: cuantos ritmos llenan un \omterm{def:g3:division:half}{medio} ¿nota?
  \item Los números que encontraste --- $1, 2, 3, 5, 8, 13, 21, 34$
        --- fueron publicados por el erudito indio Hemachandra
        alrededor del año 1150, contando exactamente esos ritmos, dos
        generaciones antes de que Fibonacci los escribiera en Italia.
        Indique su patrón definitorio en una oración y verifique
        en las últimas tres entradas de su tabla.
  \item El estilo ``Swing'' reemplaza dos octavos iguales por un
        par largo-corto: $\frac16$ entonces $\frac{1}{12}$. verificar
        que el swap es legal (misma duración total), utilizando el
        común \omterm{def:g4:fractions:def}{denominador} $12$.
  \item A \emph{trillizo} comprime tres notas iguales en una
        \omterm{def:g3:division:half}{cuarto} nota. ¿Cuánto dura cada nota? (comprueba tu respuesta).
        multiplicándolo por $3$,
        \cref{prop:g7:fractions:mult})? Misma pregunta para tres.
        notas iguales llenando un \omterm{def:g3:division:half}{medio} nota.
  \item La música de baile de los Balcanes utiliza el compás siete-ocho: cada compás
        totales $\frac78$. Verifique que el punto \omterm{def:g3:division:half}{cuarto} $+$ \omterm{def:g3:division:half}{cuarto}
        $+$ \omterm{def:g3:division:half}{cuarto} lo llena. Luego explica por qué \emph{No}
        combinación de \omterm{def:g3:division:half}{medio} notas y \omterm{def:g3:division:half}{cuarto} notas solas,
        por muchos que sean, alguna vez podrán llenar un compás siete-ocho.
        (Cuente en octavos y piense en pares y \omterm{def:g3:numbers:evenodd}{extraño}.)
\end{enumerate}
\end{problem}
